\documentclass[paper=a4, fontsize=11pt]{scrartcl}

% ── Encoding & fonts ──────────────────────────────────────────────────────────
\usepackage[T1]{fontenc}
\usepackage[utf8]{inputenc}
\usepackage{lmodern}
\usepackage{microtype}

% ── Page geometry ─────────────────────────────────────────────────────────────
\usepackage[top=1in, bottom=1in, left=1in, right=1in]{geometry}

% ── Math ──────────────────────────────────────────────────────────────────────
\usepackage{amsmath}
\usepackage{amssymb}

% ── Graphics / TikZ ───────────────────────────────────────────────────────────
\usepackage{graphicx}
\usepackage{tikz}
\usepackage{pgfplots}
\pgfplotsset{compat=1.18}
\usetikzlibrary{shapes.geometric, arrows.meta, positioning, decorations.pathreplacing, calc}

% ── Tables ────────────────────────────────────────────────────────────────────
\usepackage{booktabs}
\usepackage{multirow}

% ── Lists ─────────────────────────────────────────────────────────────────────
\usepackage{enumitem}
\setlist[itemize]{leftmargin=*, topsep=3pt, itemsep=1pt}
\setlist[enumerate]{leftmargin=*, topsep=3pt, itemsep=1pt}

% ── Hyperlinks ────────────────────────────────────────────────────────────────
\usepackage{hyperref}
\hypersetup{
  colorlinks = true,
  linkcolor  = blue!70!black,
  citecolor  = blue!70!black,
  urlcolor   = blue!70!black,
}

% ── Bibliography ──────────────────────────────────────────────────────────────
\usepackage{natbib}

% ── Header / footer (KOMA built-in) ──────────────────────────────────────────
\usepackage{scrlayer-scrpage}
\pagestyle{scrheadings}
\clearpairofpagestyles
\ohead{\normalfont\small CS 6120 · Natural Language Processing · Northeastern University}
\ihead{\normalfont\small Shourya Srivastava}
\ofoot{\pagemark}

% ── Captions ─────────────────────────────────────────────────────────────────
\usepackage{caption}
\captionsetup{font=small, labelfont=bf}

% ── Section spacing ───────────────────────────────────────────────────────────
\RedeclareSectionCommand[beforeskip=1.2ex plus .2ex, afterskip=0.6ex plus .1ex]{section}
\RedeclareSectionCommand[beforeskip=0.8ex plus .1ex, afterskip=0.4ex plus .1ex]{subsection}

% ── Title block ───────────────────────────────────────────────────────────────
\title{ArXiv Research Assistant:\\
       A Fully Local Retrieval-Augmented Generation System\\
       for Scientific Literature}
\subtitle{CS 6120 -- Natural Language Processing \\ Final Project Report}
\author{Shourya Srivastava \quad \texttt{srivastava.shou@northeastern.edu}}
\date{Spring 2026}

% ══════════════════════════════════════════════════════════════════════════════
\begin{document}
\maketitle

\noindent
\textbf{Repository:} \url{https://github.com/shouryasrivastava/arxiv-rag} \hfill
\textbf{Live demo:} \texttt{http://34.26.47.93:8501}

\medskip

% ══════════════════════════════════════════════════════════════════════════════
\section{Proposed Objective}

Large language models (LLMs) carry a hard \textbf{knowledge cutoff}: they
cannot answer questions about papers published after their training data was
collected, and they frequently hallucinate citations when asked about specific
technical claims~\citep{hallucination2023}.
The goal of this project is to build an end-to-end
\textbf{Retrieval-Augmented Generation (RAG)} system that grounds LLM responses
in over 50{,}000 peer-reviewed ArXiv papers, enabling accurate and verifiable
answers to scientific research questions without any external API calls or model
retraining.

\paragraph{Heilmeier questions.}
\begin{itemize}
  \item \textit{What are we trying to do?} Build a locally-runnable question
        answering assistant over the ArXiv scientific preprint corpus that cites
        its sources and never fabricates references.
  \item \textit{How is it done today, and what are the limits?} General LLMs
        hallucinate citations; search engines return links but cannot synthesise
        answers; existing scholarly tools (Semantic Scholar, Elicit) require
        cloud APIs.
  \item \textit{What is new in our approach?} A fully local, containerised
        stack---no API keys, reproducible on any machine with 8\,GB RAM---combining
        dense vector retrieval with a compact open-weight LLM.
  \item \textit{If successful, what difference does it make?} Researchers can
        interactively query tens of thousands of papers with grounded,
        citation-backed answers in an offline or privacy-sensitive environment.
\end{itemize}

% ══════════════════════════════════════════════════════════════════════════════
\section{Background and Motivation}

\subsection{Retrieval-Augmented Generation}
\citet{rag2020} introduced RAG for knowledge-intensive NLP tasks by coupling a
dense retriever with a seq2seq generator.
\citet{realm2020} extended retrieval to the pre-training stage, while
\citet{atlas2022} showed that a smaller RAG model can match GPT-3 on knowledge
benchmarks when retrieval quality is high.
Our system follows the \textit{modular RAG} design of~\citet{modularrag2023},
separating the retriever, context builder, and generator so each component can
be swapped independently.

\subsection{Dense Retrieval}
Sparse methods such as BM25~\citep{bm25} match on token overlap and miss
semantic synonyms.
Bi-encoder dense retrievers~\citep{dpr2020} represent queries and documents as
vectors in a shared embedding space and retrieve by maximum inner product
search.
We use \textit{all-MiniLM-L6-v2}~\citep{reimers2019sentencebert}, a
sentence-transformer distilled to 22M parameters that provides an excellent
speed--quality trade-off, producing 384-dimensional unit-norm vectors in under
25\,ms per query.

\subsection{Local LLM Inference}
The emergence of efficient open-weight models---LLaMA~\citep{llama2023},
Mistral~\citep{mistral2023}, and Llama~3~\citep{llama3}---has made high-quality
local inference practical.
Ollama~\citep{ollama2024} wraps these models in a unified REST API with 4-bit
quantisation (Q4\_K\_M), reducing the memory footprint of Llama~3.2~3B to
approximately 2\,GB and making CPU-only deployment viable.

\subsection{Data}
The Cornell-University/arxiv dataset~\citep{clement2019arxiv} on HuggingFace
contains over 1.7 million papers spanning 1991--2025.
We index the 50{,}000 most recent papers across eight CS/ML categories
(\texttt{cs.AI}, \texttt{cs.LG}, \texttt{cs.CL}, \texttt{cs.CV},
\texttt{cs.IR}, \texttt{stat.ML}, \texttt{cs.NE}, \texttt{cs.RO}),
prioritising recent work that is absent from pre-trained LLM knowledge.
Each record contributes its title and abstract (average 175 words, well within
the 512-token encoder limit) as a single indexable document.

% ══════════════════════════════════════════════════════════════════════════════
\section{Proposed Approach / Implementation Details}

\subsection{System Architecture}

The system has two phases: an offline \textbf{ingestion pipeline} that runs
once at startup, and an online \textbf{query pipeline} that runs on every user
request (Figure~\ref{fig:arch}).

\begin{figure}[h]
\centering
\begin{tikzpicture}[
  box/.style  = {rectangle, rounded corners=4pt, draw=black!50, fill=#1,
                 text centered, minimum height=0.85cm, minimum width=2.7cm,
                 font=\small},
  arr/.style  = {-{Latex[length=2mm]}, thick, color=black!65},
  darr/.style = {-{Latex[length=2mm]}, thick, dashed, color=gray!60},
  lbl/.style  = {font=\footnotesize\itshape, color=gray!75},
  node distance = 1.0cm and 1.4cm,
]
% Ingestion (left column)
\node[box=blue!10]   (hf)     {HuggingFace\\ArXiv Dataset};
\node[box=blue!10,  below=of hf]  (api)    {ArXiv API\\(fallback)};
\node[box=orange!20,below=of api] (emb)    {all-MiniLM-L6-v2\\Embedder};
\node[box=green!15, below=of emb] (chroma) {ChromaDB\\(cosine HNSW)};
\draw[arr]  (hf)  -- (emb)    node[midway,right,lbl]{stream};
\draw[darr] (api) -- (emb)    node[midway,right,lbl]{fallback};
\draw[arr]  (emb) -- (chroma) node[midway,right,lbl]{upsert};
\draw[decorate, decoration={brace,amplitude=5pt}, thick, color=blue!50]
  ([xshift=-1.5cm]hf.north west) -- ([xshift=-1.5cm]chroma.south west)
  node[midway,xshift=-0.9cm,rotate=90,anchor=south,
       font=\small\bfseries,color=blue!65]{Ingestion};

% Query (right column)
\node[box=purple!15, right=4.2cm of hf] (user)   {User Query\\(Streamlit)};
\node[box=orange!20, below=of user]     (qemb)   {Embed Query};
\node[box=green!15,  below=of qemb]     (search) {Top-$k$ Retrieval};
\node[box=red!15,    below=of search]   (llm)    {Ollama\\Llama 3.2:3b};
\node[box=purple!15, below=of llm]      (out)    {Streamed Answer\\+ Citations};
\draw[arr] (user)   -- (qemb);
\draw[arr] (qemb)   -- (search) node[midway,right,lbl]{384-d};
\draw[arr] (search) -- (llm)    node[midway,right,lbl]{context};
\draw[arr] (llm)    -- (out)    node[midway,right,lbl]{stream};
\draw[arr] (chroma.east) to[out=0,in=180] (search.west);
\draw[decorate, decoration={brace,amplitude=5pt,mirror}, thick, color=purple!50]
  ([xshift=1.5cm]user.north east) -- ([xshift=1.5cm]out.south east)
  node[midway,xshift=0.9cm,rotate=-90,anchor=south,
       font=\small\bfseries,color=purple!65]{Query};
\end{tikzpicture}
\caption{ArXiv Research Assistant architecture. The ingestion pipeline (left)
  indexes papers once; the query pipeline (right) embeds each user question,
  retrieves the top-$k$ papers, and streams a grounded answer from the local LLM.}
\label{fig:arch}
\end{figure}

\subsection{Ingestion Pipeline}
Papers are streamed from the HuggingFace \texttt{Cornell-University/arxiv}
dataset (with automatic fallback to the official ArXiv export API).
Each batch of 512 papers is embedded with \textit{all-MiniLM-L6-v2} and
immediately upserted into ChromaDB~\citep{chromadb2023}, so the index is
queryable from the very first batch.
The pipeline supports \textbf{incremental resume}: on restart it skips
already-indexed document IDs.
The ChromaDB collection uses a cosine HNSW index for sub-millisecond
approximate nearest-neighbour search at 50k-document scale.

\paragraph{Preprocessing.}
Each document is the concatenation of the paper's title and abstract (mean
175 words, $\sigma \approx 45$), fitting within the 512-token encoder limit.
No additional chunking is required at the abstract level.

\subsection{Query Pipeline}
\begin{enumerate}
  \item \textbf{Query embedding.} The user question is encoded by
        \textit{all-MiniLM-L6-v2} (unit-norm, 384-d, $<$25\,ms).
  \item \textbf{Dense retrieval.} ChromaDB returns the top-$k$ papers by
        cosine similarity ($k = 5$ by default; adjustable via UI slider).
        An optional category filter (e.g.\ \texttt{cs.CV}) can be applied
        via a metadata predicate.
  \item \textbf{Prompt construction.} A structured prompt is built with a
        strict system instruction, numbered abstract blocks with
        title/authors/date/URL, and the user question.
  \item \textbf{Streaming generation.} The prompt is sent to Ollama's REST
        endpoint with \texttt{stream: true}; tokens are displayed live in
        Streamlit.
  \item \textbf{Citation cards.} Each retrieved paper is shown with its title,
        relevance score, categories, date, and a clickable
        \texttt{arxiv.org/abs/<id>} link.
\end{enumerate}

\subsection{Hallucination Mitigation}
The system prompt explicitly constrains the model:
\begin{quote}\small
\textit{``Answer using ONLY the information provided in the retrieved abstracts.
If the context does not contain enough information, say so explicitly.
Do NOT speculate beyond the provided context.''}
\end{quote}
Generation uses temperature 0.1 and top-$p$ = 0.9 to minimise deviation from
source material. Across our 30-query evaluation the system produced
\textbf{zero hallucinated citations}.

\subsection{Evaluation}

\begin{table}[h]
\centering
\caption{Retrieval quality on 30 manually curated research queries.}
\label{tab:retrieval}
\begin{tabular}{lcccc}
\toprule
\textbf{Method} & \textbf{P@1} & \textbf{P@3} & \textbf{P@5} & \textbf{MRR} \\
\midrule
BM25 (baseline)          & 0.53 & 0.61 & 0.67 & 0.58 \\
all-MiniLM-L6-v2 (ours)  & \textbf{0.77} & \textbf{0.79} & \textbf{0.81} & \textbf{0.79} \\
\bottomrule
\end{tabular}
\end{table}

Dense retrieval outperforms BM25 by +24\,pp at P@1.
End-to-end latency on a CPU-only GCP \texttt{n2-standard-4} instance is
$\approx$51\,s (embedding: 22\,ms, retrieval: 8\,ms, LLM generation:
$\approx$51\,s for 200 tokens).

\subsection{Deployment}
The full stack is containerised with Docker Compose into two services:
\texttt{ollama} (port 11434) and \texttt{app} (Python 3.11-slim, port 8501).
Docker volumes persist ChromaDB and Ollama model weights across restarts.
A startup script pre-warms the LLM with a single-token dummy request so
weights are in RAM before the first user query.
The system is deployed on GCP at \texttt{http://34.26.47.93:8501}.

% ══════════════════════════════════════════════════════════════════════════════
\section{Submission Guidelines}

\begin{itemize}
  \item \textbf{Repository (public):}
        \url{https://github.com/shouryasrivastava/arxiv-rag} ---
        contains \texttt{src/}, \texttt{scripts/}, \texttt{Dockerfile},
        \texttt{docker-compose.yml}, \texttt{requirements.txt}, and this
        write-up under \texttt{writeup/}.
  \item \textbf{Live demo endpoint:}
        \texttt{http://34.26.47.93:8501} --- publicly accessible Streamlit UI
        running on a GCP \texttt{n2-standard-4} VM.
  \item \textbf{One-command local setup:}
        \texttt{docker compose up -{-}build} (requires Docker Desktop and
        $\geq$8\,GB free RAM; model download $\approx$2\,GB on first run).
  \item \textbf{Author:}
        Shourya Srivastava --- system architecture, ingestion pipeline, RAG
        pipeline, Streamlit frontend, Docker deployment, GCP setup, evaluation.
\end{itemize}

% ── References ────────────────────────────────────────────────────────────────
\bibliographystyle{abbrvnat}
\bibliography{references}

\end{document}
